% Goal Module for drawing grid scenarios
% Required packages (should be in main.tex preamble):
% \usepackage{xparse}
% \usepackage{tikz}
% \usetikzlibrary{positioning}
% \usepackage{xstring}

% Color definitions (if not already defined)
\definecolor{colorF}{HTML}{3B82F6} % blue!45 approx
\definecolor{colorL}{HTML}{EF4444} % red!60 approx
\definecolor{colorR}{HTML}{22C55E} % Green!60
\definecolor{colorO}{HTML}{F97316} % orange!90

% The macro definition
\NewDocumentCommand{\drawGridWithHighlightsAndBox}{mmmmmmm}{
    \begin{tikzpicture}
        \def\nodesize{0.4cm}
        \def\spacing{0.6cm}

        \pgfmathsetmacro{\maxX}{int(#1)}
        \pgfmathsetmacro{\maxY}{int(#2)}

        % Define coordinates
        \foreach \x in {0,...,\maxX} {
            \foreach \y in {0,...,\maxY} {
                \coordinate (n\x\y) at (\x*\spacing, \y*\spacing);
            }
        }

        % Draw grid lines
        \foreach \x in {0,...,\numexpr\maxX-1} {
            \foreach \y in {0,...,\maxY} {
                \pgfmathtruncatemacro{\xp}{\x+1}
                \draw[gray!40, thin] (n\x\y) -- (n\xp\y);
            }
        }
        \foreach \x in {0,...,\maxX} {
            \foreach \y in {0,...,\numexpr\maxY-1} {
                \pgfmathtruncatemacro{\yp}{\y+1}
                \draw[gray!40, thin] (n\x\y) -- (n\x\yp);
            }
        }

        % Draw default nodes (background circles)
        \foreach \x in {0,...,\maxX} {
            \foreach \y in {0,...,\maxY} {
                \node[circle, draw=gray!20, thin, fill=white, minimum size=\nodesize, inner sep=0pt]
                    at (n\x\y) {};
            }
        }

        % Highlighted nodes with optional labels
        % Format: color/x/y/label
        \foreach \entry in {#3} {
            \StrCut{\entry}{/}{\color}{\rest}
            \StrCut{\rest}{/}{\xval}{\rest}
            \StrCut{\rest}{/}{\yval}{\label}
            \IfStrEq{\label}{}{\edef\labelText{}}{\edef\labelText{\label}}
            \node[circle, draw=black, thick, fill=\color, minimum size=\nodesize, inner sep=0pt, font=\tiny\bfseries]
                at (n\xval\yval) {\labelText};
        }

        % Draw bounding box
        \draw[thick, accent, dashed, rounded corners=2pt]
            ([shift={(-0.35cm,-0.35cm)}]n#4#5)
            rectangle
            ([shift={(0.35cm,0.35cm)}]n#6#7);
    \end{tikzpicture}
}
